\section{Concluding remarks}
Using historical time series of implied volatility surfaces for the S\&P 500 and the Euro Stoxx 50, we have shown empirically that a large part of the variability of the at-the-money-forward (ATM) implied volatility for time-to-maturities ranging from 1 month to 24 months can be explained by two features, namely the weighted average of the underlying asset past returns and the weighted average of the past squared returns. As the time-to-maturity increases, the part of variability explained by these two features decreases but remains important. Thus, our empirical study allows to extend the one of \cite{guyon2022volatility}  who focused on implied volatility indices and realized volatility instead of vanilla options implied volatilities. In Section \ref{sec:ssvi}, we have then introduced a parsimonious version (see Definition \ref{def:parsimonious_ssvi}) of the SSVI parameterization of \cite{gatheral2014arbitrage} that depends on four parameters only ($a$, $p$, $\rho$ and $\eta$) and that still achieves a reasonable fit to the implied volatility surfaces in our two data sets. This parsimonious version is essentially obtained by considering a parametric form of the ATM total variance thus avoiding to have one parameter per time-to-maturity. In the last section, we demonstrate that the variations of the two parameters $a$ and $p$ ruling the ATM implied volatility in the parsimonious SSVI parameterization can also be widely explained by the two features that we mentioned earlier. Based on this observation, we introduce a new model for the joint dynamics of the underlying asset price and the full implied volatility surface (there is no restriction in the range of time-to-maturities and strikes that one wants to project) embedding the path-dependency of the implied volatility with respect to the underlying price. On the one hand, the underlying asset price is modelled using the path-dependent volatility model of \cite{guyon2022volatility}. On the other hand, the residuals of the parameters $a$ and $p$ are modelled through Ornstein-Uhlenbeck processes and the parameters $\rho$ and $\eta$ are modelled as Jacobi processes which are guaranteeing the absence of static arbitrage in the simulated surfaces.  Extensive details on how to calibrate this new model are provided. Finally, we show the high consistency of this model with historical data. The study of impact of this new model for applications in asset management, risk management and hedging is left for future research. 
